\section{Desenvolvimento}
\label{sec:desenvolvimento}

Esta seção descreve como o estudo foi concretamente implementado, desde a organização da ferramenta experimental até a execução dos experimentos e a preparação dos dados que embasam a análise apresentada na seção de Resultados e Discussão.

\subsection{Arquitetura da Ferramenta Experimental}

A implementação foi estruturada como uma pequena \emph{ferramenta de experimentação} em \texttt{Python}, organizada em módulos com responsabilidades bem definidas. Essa separação favorece a reprodutibilidade do estudo e permite que novos cenários de teste sejam adicionados com o mínimo de alterações no código existente.

De forma geral, a arquitetura é composta pelos seguintes blocos lógicos:

\begin{itemize}
    \item \textbf{Núcleo do algoritmo genético}:
    responsável pelo laço evolutivo principal, inicialização da população, aplicação dos operadores genéticos e atualização da geração corrente.
    \item \textbf{Módulo de representação e aptidão}:
    encapsula a codificação em forma de permutação para o problema das \textit{n}-rainhas e a função de avaliação baseada na contagem de pares de rainhas em conflito diagonal, conforme definido na Seção~\ref{sec:fundamentacao}.
    \item \textbf{Módulo de operadores genéticos}:
    implementa seleção, cruzamento e mutação de forma independente do problema, recebendo apenas indivíduos e retornando descendentes.
    \item \textbf{Módulo de controle populacional}:
    concentra a lógica das quatro estratégias de substituição comparadas neste trabalho (população fixa, substituição geracional, elitismo e controle por idade/variável).
    \item \textbf{Módulo de experimentos}:
    responsável por ler as configurações de cada cenário (tamanho do tabuleiro, limites de gerações e tempo, número de execuções independentes, entre outros), orquestrar a execução repetida do algoritmo genético e registrar as métricas de interesse em arquivos de saída.
    \item \textbf{Módulo de registro e persistência}:
    implementa a escrita de logs e a serialização dos resultados em formato tabular (\texttt{.csv}), de modo a facilitar o pós-processamento estatístico.
\end{itemize}

Essa organização modular permite que o mesmo núcleo de algoritmo genético seja utilizado com diferentes políticas de controle populacional, modificando-se apenas a configuração carregada pelo módulo de experimentos. Além disso, a presença de um módulo dedicado a registro de resultados simplifica a automatização de baterias extensas de testes, uma vez que a ferramenta pode ser executada em modo não interativo a partir de linha de comando.

\subsection{Implementação dos Operadores Genéticos}

A implementação segue as decisões de projeto descritas na Seção~\ref{sec:metodologia}. Cada indivíduo é representado por um vetor de inteiros de tamanho $n$, onde a posição do vetor corresponde à coluna do tabuleiro e o valor armazenado indica a linha ocupada pela rainha naquela coluna. Como essa representação garante que não haja conflitos em linhas e colunas, a função de aptidão concentra-se em contabilizar apenas conflitos diagonais.

A função de aptidão percorre pares de colunas e verifica se a diferença de índices é igual à diferença absoluta entre as linhas correspondentes. Quando essa condição se verifica, o par de rainhas está em conflito diagonal e uma penalização é aplicada. Para fins de comparação entre estratégias de controle populacional, o valor de aptidão é calculado de forma consistente em todos os cenários, e uma solução ótima é caracterizada pela ausência de conflitos diagonais.

A seleção é realizada por torneio binário: dois indivíduos são escolhidos aleatoriamente da população e o de melhor aptidão é selecionado como pai. Esse procedimento é repetido até que seja formada a quantidade necessária de pais para gerar a prole da próxima geração. Esse método é simples, eficiente e não requer normalização explícita dos valores de aptidão.

O cruzamento adotado é o \textit{Order Crossover} (OX), adequado a representações permutacionais. Na implementação, dois pontos de corte são sorteados aleatoriamente ao longo do vetor. O segmento interno entre os cortes é copiado diretamente de um dos pais para o descendente, preservando a ordem relativa dos elementos. Em seguida, as posições restantes são preenchidas com os genes do outro pai, respeitando a ordem em que aparecem e evitando duplicações. O resultado é um descendente que herda blocos estruturais de ambos os pais, mantendo a propriedade de permutação válida.

A mutação é implementada na forma de troca de posições (\emph{swap mutation}). Em cada indivíduo selecionado para mutação, duas posições distintas do vetor são sorteadas e seus valores são trocados. Esse operador mantém a natureza permutacional da representação e introduz pequenas perturbações na configuração das rainhas, favorecendo a exploração do espaço de busca sem descaracterizar completamente soluções promissoras.

Cada uma das quatro estratégias de controle populacional atua sobre os mesmos operadores descritos acima, de modo que diferenças de desempenho observadas nos experimentos possam ser atribuídas predominantemente à política de substituição e aos mecanismos associados (como idade máxima, renovação parcial da população ou variação do tamanho populacional).

\subsection{Execução dos Experimentos}

A ferramenta experimental foi configurada para executar automaticamente todas as combinações de parâmetros definidos na Seção~\ref{sec:metodologia}. Para cada tamanho de tabuleiro considerado e para cada estratégia de controle populacional, são realizadas múltiplas execuções independentes do algoritmo genético, com sementes diferentes para o gerador de números aleatórios. Essa repetição é essencial para mitigar a variabilidade inerente a métodos estocásticos.

Cada execução segue o fluxo descrito a seguir:

\begin{enumerate}
    \item \textbf{Inicialização}: a partir da configuração fornecida, a população inicial é gerada por embaralhamento aleatório de sequências de $0$ a $n-1$. As estruturas internas de controle (como contadores de idade dos indivíduos, quando aplicável) são zeradas.
    \item \textbf{Laço evolutivo}: em cada geração, a função de aptidão é avaliada para todos os indivíduos (ou para aqueles afetados por modificações, dependendo da estratégia). Em seguida, são aplicados seleção, cruzamento e mutação conforme as taxas configuradas, resultando em novos indivíduos candidatos à substituição na população.
    \item \textbf{Aplicação da estratégia de controle populacional}: de acordo com o esquema em avaliação, a população é atualizada. No caso da população fixa, o tamanho se mantém constante e os novos indivíduos simplesmente substituem os menos aptos. Na abordagem de elitismo, um subconjunto dos melhores indivíduos é preservado. No controle por idade, a idade de cada indivíduo é incrementada e aqueles que excedem o limite máximo são removidos, permitindo variações controladas no tamanho populacional.
    \item \textbf{Verificação dos critérios de parada}: após cada geração, verifica-se se algum dos critérios de parada foi satisfeito: (i) solução ótima encontrada; (ii) limite máximo de gerações atingido; ou (iii) limite de tempo excedido. Em qualquer desses casos, a execução é encerrada.
    \item \textbf{Registro das métricas}: ao final de cada execução, são registrados o tempo total decorrido, o número de gerações percorridas, o critério de parada acionado e um indicador de sucesso, além de informações de configuração (estratégia de controle adotada, tamanho do tabuleiro, semente aleatória, entre outras).
\end{enumerate}

Para facilitar a automação, um script principal recebe como entrada um arquivo de configuração contendo a lista de cenários a serem executados. Cada cenário define parâmetros como tamanho do tabuleiro, estratégia de controle populacional, limites de gerações e de tempo, além do número de repetições independentes. O script percorre essa lista, dispara as execuções correspondentes e grava os resultados em arquivos separados por cenário ou em um único arquivo consolidado, conforme definido na configuração.

\subsection{Organização dos Dados e Pré-processamento}

Os dados produzidos pela ferramenta experimental são armazenados em arquivos no formato \texttt{.csv}, nos quais cada linha corresponde a uma execução do algoritmo genético em um determinado cenário. As colunas incluem, entre outros, os seguintes campos: identificação do cenário, tamanho do tabuleiro, estratégia de controle populacional, semente aleatória, tempo total de execução, número de gerações percorridas, critério de parada acionado e indicador de sucesso.

Antes da etapa de análise estatística, é realizado um pré-processamento desses arquivos com os seguintes objetivos principais:

\begin{itemize}
    \item \textbf{Validação e limpeza}: verificação de registros incompletos ou inconsistentes (por exemplo, execuções interrompidas por falhas externas) e remoção ou correção desses casos, quando necessário;
    \item \textbf{Cálculo de métricas derivadas}: criação de colunas adicionais a partir das métricas brutas, como taxas de sucesso por cenário, tempo médio até a solução, número médio de gerações por execução bem-sucedida e proporção de cada critério de parada acionado;
    \item \textbf{Agrupamento por estratégia e tamanho de problema}: reorganização dos dados em estruturas que facilitem a comparação entre estratégias de controle populacional, mantendo separadas as diferentes dificuldades (valores de $n$) e consolidando resultados de execuções independentes.
\end{itemize}

O pré-processamento é implementado em \texttt{Python}, utilizando bibliotecas de manipulação de dados para carregar os arquivos \texttt{.csv}, aplicar filtros e agregar estatísticas. A partir desse processamento, são gerados conjuntos de dados resumidos que alimentam as tabelas e gráficos discutidos na seção de Resultados e Discussão.

Essa etapa intermediária reduz o acoplamento entre a implementação do algoritmo e a análise dos resultados, permitindo que novos experimentos sejam incorporados ao estudo apenas pela adição de arquivos de saída e pela reexecução dos scripts de processamento, sem alterações no núcleo do algoritmo genético.
